% PDE 论文通用模板 —— 预置偏微分方程写作常用宏与定理环境
% 引擎：pdflatex（中文请切换到 ctex 模板 + xelatex）
\documentclass[11pt,a4paper]{article}

\usepackage{amsmath,amssymb,amsthm,mathtools}
\usepackage{geometry}
\usepackage{graphicx}
\usepackage{hyperref}
\geometry{margin=2.7cm}
\hypersetup{colorlinks=true,linkcolor=blue,citecolor=blue}

% ---------- 定理环境 ----------
\newtheorem{theorem}{Theorem}[section]
\newtheorem{lemma}[theorem]{Lemma}
\newtheorem{proposition}[theorem]{Proposition}
\newtheorem{corollary}[theorem]{Corollary}
\theoremstyle{definition}
\newtheorem{definition}[theorem]{Definition}
\newtheorem{assumption}[theorem]{Assumption}
\theoremstyle{remark}
\newtheorem{remark}[theorem]{Remark}

% ---------- PDE 常用宏 ----------
\newcommand{\R}{\mathbb{R}}
\newcommand{\N}{\mathbb{N}}
\newcommand{\Om}{\Omega}
\newcommand{\paOm}{\partial\Omega}
\newcommand{\eps}{\varepsilon}
\newcommand{\weakto}{\rightharpoonup}
\newcommand{\norm}[1]{\left\lVert #1\right\rVert}
\newcommand{\abs}[1]{\left\lvert #1\right\rvert}
\newcommand{\inner}[2]{\left\langle #1,#2\right\rangle}
\newcommand{\dd}{\,\mathrm{d}}
\DeclareMathOperator{\divg}{div}
\DeclareMathOperator{\curl}{curl}
\DeclareMathOperator{\supp}{supp}
\DeclareMathOperator{\dist}{dist}
\newcommand{\pde}[1]{\eqref{#1}}

\title{Well-posedness and Long-Time Dynamics of a Nonlinear Parabolic Equation}
\author{Yang Jianwei}
\date{\today}

\begin{document}
\maketitle

\begin{abstract}
We establish the global well-posedness of a nonlinear parabolic equation
in $H^1_0(\Omega)$ and prove the existence of a global attractor with
finite fractal dimension. The key ingredient is a uniform
$H^2$ estimate obtained by an Alikakos--Moser iteration.
\end{abstract}

\section{Introduction}

Let $\Om\subset\R^n$ be a bounded domain with smooth boundary $\paOm$.
We study the initial--boundary value problem
\begin{equation}\label{P}
  \begin{cases}
    u_t - \divg\bigl(a(x)\nabla u\bigr) + f(u) = g(x), & (x,t)\in\Om\times\R^+,\\[2pt]
    u = 0, & (x,t)\in\paOm\times\R^+,\\[2pt]
    u(x,0)=u_0(x), & x\in\Om,
  \end{cases}
\end{equation}
where $a\in C^1(\overline\Om)$ is uniformly elliptic, i.e.\
$\lambda|\xi|^2\le a(x)\xi\cdot\xi\le\Lambda|\xi|^2$ for all
$\xi\in\R^n$.

\begin{assumption}\label{A1}
The nonlinearity $f\in C^1(\R)$ satisfies
$f'(s)\ge -c_0$ and $|f(s)|\le c_1(1+|s|^{p-1})$ with
$1\le p\le \frac{n+2}{n-2}$.
\end{assumption}

\begin{theorem}[Well-posedness]\label{T1}
Under Assumption~\ref{A1}, for every $u_0\in L^2(\Om)$ problem~\pde{P}
admits a unique weak solution
$u\in C([0,\infty);L^2(\Om))\cap L^2_{\mathrm{loc}}(0,\infty;H^1_0(\Om))$.
Moreover, the solution depends continuously on the initial data.
\end{theorem}

\section{A priori estimates}

\begin{lemma}[Energy identity]\label{L1}
For any $t>0$,
\begin{equation}
  \frac12\frac{\dd}{\dd t}\norm{u(t)}_{L^2}^2
  + \int_{\Om}a(x)|\nabla u|^2\dd x
  + \int_{\Om}f(u)u\dd x
  = \int_{\Om}gu\dd x .
\end{equation}
\end{lemma}

\begin{proof}
Multiply \pde{P} by $u$, integrate over $\Om$, and use
$\inner{-\divg(a\nabla u),u}=\int_\Om a|\nabla u|^2\dd x$.
\end{proof}

\begin{remark}
By Assumption~\ref{A1} and Poincaré's inequality, Lemma~\ref{L1} yields
a uniform bound $\norm{u(t)}_{L^2}\le C(1+\norm{u_0}_{L^2}e^{-\mu t})$,
which gives the existence of an absorbing set in $L^2(\Om)$.
\end{remark}

\section{Proof of Theorem~\ref{T1}}

\begin{proof}
We use a Faedo--Galerkin approximation. Let $\{w_k\}_{k\ge1}$ be the
eigenfunctions of $-\divg(a\nabla\cdot)$ with Dirichlet boundary
condition. For each $m$, let $u_m=\sum_{k=1}^m c_{mk}(t)w_k$ solve the
projected ODE system. Lemma~\ref{L1} provides uniform bounds, and
compactness (Aubin--Lions) allows passing to the limit $m\to\infty$.
Uniqueness follows from the monotonicity of $f$ and Gronwall's inequality.
\end{proof}

For the classical background and related results, see~\cite{Temam,Lady}.

\begin{thebibliography}{9}
\bibitem{Temam} R.~Temam, \emph{Infinite-Dimensional Dynamical Systems in
Mechanics and Physics}, 2nd ed., Springer, 1997.
\bibitem{Lady} O.~A.~Ladyzhenskaya, V.~A.~Solonnikov and N.~N.~Ural'ceva,
\emph{Linear and Quasi-linear Equations of Parabolic Type}, AMS, 1968.
\end{thebibliography}

\end{document}
